%!TeX program = pdflatex
% Russian CV for Legal Audience - One Page
\documentclass[10pt,a4paper]{article}

\usepackage[T2A]{fontenc}
\usepackage[utf8]{inputenc}
\usepackage[russian]{babel}
\usepackage{microtype}
\usepackage{hyperref}
\usepackage{geometry}
\usepackage{enumitem}
\usepackage{titlesec}

% Имя
\newcommand{\myname}{Алишер Джузгенбаев}

% Шрифты
\usepackage{libertine}

% Определение шрифта для заголовка
\newcommand{\namefont}[1]{{\normalfont\bfseries\Large{#1}}}

% Настройка заголовков секций
\titleformat{\section}
  {\normalfont\bfseries}
  {}
  {0em}
  {}
  [\titlerule]
  
\titlespacing*{\section}
  {0pt}{10pt}{5pt}

% Настройка полей для одной страницы
\geometry{
    a4paper,
    left=1.5cm,
    right=1.5cm,
    top=1.5cm,
    bottom=1.5cm
}

% Убираем отступы параграфов
\setlength\parindent{0em}

% Компактные списки
\setlist[itemize]{leftmargin=1em, topsep=0pt, itemsep=1pt, parsep=0pt}

% Настройка гиперссылок
\hypersetup{
    colorlinks  = true,
    urlcolor    = blue,
    linkcolor   = black,
    pdfauthor   = {\myname},
    pdftitle    = {\myname: Резюме},
}

\begin{document}
\pagestyle{empty} % Убираем номера страниц

% Заголовок
\begin{center}
    \namefont{\myname}
    
    \vspace{0.3em}
    Northwestern University • Chicago, IL, USA \\
    \href{mailto:juzgenbayev@u.northwestern.edu}{juzgenbayev@u.northwestern.edu} • 
    +7 (747) 651-5632 • 
    \href{https://alisher.site}{alisher.site}
\end{center}

\section{Образование}

\begin{itemize}
    \item[] \textbf{J.D.} (Доктор права), Northwestern Pritzker School of Law, 2027 (ожидается)
    \item[] \textbf{Ph.D.} (Политология), Northwestern University, 2027 (ожидается)
    \item[] \textbf{M.A.} (Политология), Northwestern University, 2022
    \item[] \textbf{B.A.} (Политология и международные отношения), \textit{summa cum laude}, Назарбаев Университет, 2020
\end{itemize}

\section{Юридический опыт}

\begin{itemize}
    \item[] \textbf{Skadden, Arps, Slate, Meagher \& Flom LLP} — Летний ассоциат (2025)
    \item[] \textbf{Legal Aid Chicago} — Стажёр, группа защиты прав потребителей (2024)
    \item[] \textbf{Northwestern University Law Review} — Старший редактор по эмпирическим исследованиям, Том 119 (2023--24)
    \item[] \textbf{Федеральный окружной суд США (Северный округ Иллинойса)} — Судебный стажёр (2023)
\end{itemize}

\section{Научные публикации}

\textbf{Статьи в рецензируемых журналах:}
\begin{itemize}
    \item Трошев А., Джузгенбаев А. <<Расходящиеся траектории постсоветских конституционных судов>> [EN]. \textit{Constitutional Studies}, 11(1), 2025, pp. 271--315.
    \item Джузгенбаев А. <<Фрейминг судебной власти: эффекты партийных, процедурных и популистских фреймов на восприятие высших судов в Чехии>> [EN]. \textit{Journal of Law and Courts}, 13(1), 2025, pp. 97--121.
\end{itemize}

\textbf{Главы в книгах и статьи в юридических обозрениях:}
\begin{itemize}
    \item Джузгенбаев А. <<Путь в никуда: Ограничение права на иск и административная дискреция как угроза защите прав потребителей>> [EN]. \textit{Northwestern University Law Review}, 120, 2026 (в печати).
    \item Трошев А., Джузгенбаев А. <<Инструментализация конституционного права в Центральной Азии>> [EN] в \textit{Research Handbook on Law and Political Systems}. Edward Elgar, 2023, pp. 139--168.
\end{itemize}

\section{Области исследований}

\textbf{Диссертация:} <<Административная юстиция и правосознание в постсоветских странах>> [EN] \\
\textbf{Специализация:} Сравнительное конституционное право • Судебные системы постсоветского пространства • Эмпирические правовые исследования • Административное право

\section{Гранты и награды}

\begin{itemize}
    \item Президентская стипендия, Northwestern University (2025)
    \item Грант для международных диссертационных исследований, Институт Баффетта (2025)
    \item Грант для аспирантских исследований, Northwestern University (2023)
    \item Полуфиналист, Чемпионат Европы по дебатам среди университетов (2020)
    \item Программа для иностранных студентов, Университет Висконсин-Мэдисон (2019)
\end{itemize}

\section{Конференции и доклады}

\begin{itemize}
    \item <<Административная юстиция и правосознание в Казахстане>> [EN] — Ежегодная конвенция Law \& Society, Чикаго (2025)
    \item <<Влияние громких дел на общественное восприятие домашнего насилия: динамика идентичности в Казахстане>> [EN] — Семинар по моделям и экспериментам, Университет Висконсин-Мэдисон (2024)
    \item <<Популистские фреймы и легитимность судебной власти: экспериментальное исследование в Чешской Республике>> [EN] — Генеральная конференция ECPR, Прага (2024)
    \item <<Клиентелизм, конституционные суды и высокая политика в Центральной Азии>> [EN] — Конвенция ASEEES, Чикаго (2022)
\end{itemize}

\section{Преподавание}

\textbf{Ассистент преподавателя:} Гражданский процесс • Вещное право • Сравнительная политика • Международное право • Право, политика и общество • Международные отношения (2021--2025)

\section{Профессиональные навыки}

\begin{itemize}
    \item[] \textbf{Языки:} Русский (родной) • Английский (свободно) • Украинский (чтение) • Казахский (средний)
    \item[] \textbf{Технические навыки:} R, Python, \LaTeX, Qualtrics, эмпирический анализ данных
    \item[] \textbf{Членство:} Law \& Society Association • Association for Slavic, East European, and Eurasian Studies • American Political Science Association
\end{itemize}

\end{document}